\label{sec:const}
%==================================================================

\begin{exm}[простой пучок]\label{exm:const}
Пусть $A$ --- абелева группа, а предпучок $\Shv{F}$ определён формулой
$\Shv{F}(U)=A$ для всех непустых открытых $U$ (и
$\Shv{F}(\varnothing)=0$), все ограничения $\rho$ --- тождественные.
Тогда:
\begin{bullets}
\item $\Shv{F}_x=A$ для всякой $x$ (предел по окрестностям ничего не
 меняет, ведь все $\rho$ --- тождества), поэтому
 $\et(\Shv{F})=X\times A$, и $\pi\colon X\times A\to X$ --- обычная
 проекция;
\item топология на $X\times A$ --- произведение топологии $X$ и
  дискретной топологии на $A$ (это как раз наиболее тонкая топология,
 при которой все <<константные>> сечения непрерывны);
\item $\Shv{F}^{\#}(U)$ --- непрерывные отображения $U\to A$, т.\,е.
 \emph{локально постоянные} функции, а на связном $U$ --- просто
 константы.
\end{bullets}
Проверка (F1) для этого примера --- обычное утверждение из топологии:
локально постоянные функции на связном множестве постоянны, а на
произвольном --- определяются значениями на связных компонентах.
\end{exm}

\begin{risunok}{Постоянный предпучок: $\et(\Shv{F})=X\times A$ с
 дискретным $A$; сечение --- локально постоянная функция $U\to A$.
  \label{fig:const}}
\begin{tikzpicture}[>=Stealth,font=\small]
  % база X
  \draw[fig base] (0,0) ellipse (3.2 and 0.7);
  \node[below] at (0,-0.85) {$X$};
  % уровни {x} x {a}
  \foreach \y in {1.5,2.5,3.5} \draw[fig fiber] (-3.2,\y) -- (3.2,\y);
  \foreach \y/\a in {1.5/{a_1},2.5/{a_2},3.5/{a_3}} {
    \node[anchor=east,font=\footnotesize] at (-3.35,\y)
      {$\{x\}\times\{\a\}$};
    \foreach \x in {-2.6,-1.6,-0.6,0.4,1.4,2.4}
      \fill (\x,\y) circle (2.2pt);
  }
  % волокна
  \foreach \x/\top in {-2.6/0.408,-1.6/0.606,-0.6/0.687,
                       0.4/0.694,1.4/0.629,2.4/0.463}
    \draw[fig fiber] (\x,\top) -- (\x,3.5);
  % открытые U_1, U_2
  \begin{scope}
    \clip (0,0) ellipse (3.2 and 0.7);
    \fill[primary!10] (-2.1,0) ellipse (1.45 and 0.65);
    \fill[success!10] (1.5,0) ellipse (1.45 and 0.65);
  \end{scope}
  \draw[fig open] (-2.1,0) ellipse (1.45 and 0.65);
  \draw[fig open2] (1.5,0) ellipse (1.45 and 0.65);
  \node[primary] at (-2.1,-1.1) {$U_1$};
  \node[success] at (1.5,-1.1) {$U_2$};
  % локально постоянные сечения
  \draw[fig sect] (-3.15,2.5) -- (-0.7,2.5);
  \draw[fig sect] (0.1,1.5) -- (2.9,1.5);
  \node[danger] at (-1.9,2.68) {$a_2$};
  \node[danger] at (1.5,1.68) {$a_1$};
  % проекция
  \draw[->,thick] (3.75,3.5) -- (3.75,0.5);
  \node at (3.95,2.0) {$\pi$};
  \node[anchor=south east] at (-3.3,3.85) {$X\times A=\et(\Shv{F})$};
\end{tikzpicture}
\end{risunok}

Пояснение к рисунку~\ref{fig:const}. Овал внизу --- база $X$;
всё вместе --- $\et(\Shv{F})=X\times A$, где $A$ дискретно,
поэтому <<этажёрка>> состоит из отдельных горизонтальных
рядов точек:
горизонтальные пунктирные линии --- уровни $\{x\}\times\{a\}$
для трёх значений $a_1,a_2,a_3$. Синее и зелёное --- открытые
$U_1,U_2\subset X$. Красные отрезки --- сечение: над $U_1$ оно
постоянно равно $a_2$ (лежит на одном уровне), над $U_2$ ---
равно $a_1$: именно локально постоянная функция $U\to A$.
Стрелка $\pi$ --- проекция на базу.

\begin{bookbox}
\booksrc{Годеман, гл.~II, \S\,1, п.~1.4 <<Простые пучки>>, с.~134:
 простой пучок с базой $X$ и слоем $A$ --- пучок, порождённый предпучком
 $U\mapsto A$ с тождественными ограничениями;
 $\et=X\times A$ --- произведение топологии $X$ и дискретной топологии
 на $A$; сечения над $U$ канонически отождествляются с локально
 постоянными отображениями $U\to A$.
}
\end{bookbox}

%==================================================================
